% Source: https://artofproblemsolving.com/community/c3h1591098_question_on_statistics
% Stable identity: aops-topic-1591098-q01
One counterexample is enough. Take the data set
\[
1,\ 1,\ 2,\ 3,\ 8 .
\]
Its mode is $1$ (the only value that appears twice), its median is the middle value $2$, and its mean is $\dfrac{1+1+2+3+8}{5}=3$. Then
\[
3\cdot\text{median}-2\cdot\text{mean}=6-6=0\ne1=\text{mode},
\]
so the formula fails. $\blacksquare$

The relation cannot be proved because it is not a theorem: it is an \emph{empirical rule of thumb} (due to Karl Pearson) that holds approximately for many moderately skewed, single-peaked distributions. It is exact in special cases, for example for symmetric data with a single mode, where mode, median and mean coincide.
